
\subsection{光子场}
\topictitle{螺旋度与表示的对应}
由 3.4 节的核心结果 (3.93), $(A,B)$ 表示中零质量粒子的螺旋度为 $\lambda = B - A$。
对于自旋-1 粒子 ($\lambda = \pm 1$), 有两种最小的不可约表示：
\begin{center}
\begin{tabular}{c c c}
\hline
表示 & $\lambda$ & 物理含义 \\
\hline
$(1,0)$ & $0-1=-1$ & 反自对偶 (anti-self-dual) \\
$(0,1)$ & $1-0=+1$ & 自对偶 (self-dual) \\
\hline
\end{tabular}
\end{center}
由 3.2 节关于空间反演的讨论, 宇称变换将 
$$\mathbb{P}:(A,B) \to (B,A)\to\mathbb{P}: (1,0) \leftrightarrow (0,1)$$
因此, 为了保持空间反演对称性, 必须取直和表示
\begin{equation}
    (1,0) \oplus (0,1)
\end{equation}
这是一个 $3+3=6$ 维表示, 恰好对应四维反对称张量 $f_{\mu\nu} = -f_{\nu\mu}$ 的 6 个独立分量。
在标准动量 $k^\mu=(\kappa,0,0,\kappa)$ 处, 系数函数由 ISO(2) 约束和 CG 系数完全确定,
结果极其简单——仅一个分量非零, 值为 $1/\sqrt{2|\boldsymbol{p}|}$。
推广到任意动量, 一般公式给出
\begin{equation}
    v_{ab}(\mathbf{k},\lambda) = \sqrt{\frac{1}{2k^0}}
    \sum_{\bar{a}\bar{b}}
    \bigl[R(\hat{k})\bigr]_{a\bar{a}}\,
    \bigl[R(\hat{k})\bigr]_{b\bar{b}}\,
    \langle A,-A;\,B,+B\,|\,j,\lambda\rangle
    \tag{3.108}
\end{equation}
其中 $R(\hat{k})$ 是将 $\hat{z}$ 旋转到 $\hat{k}$ 方向的旋转矩阵,
作用在球面基指标 $a$ 和 $b$ 上。
最重要的问题是：如何将 (3.108) 中的指标 $(a,b)$ 转换到时空指标 $\mu\nu$ 上。


\paragraph*{自对偶与反自对偶分解。}
四维反对称张量 $f_{\mu\nu}=-f_{\nu\mu}$ 有 6 个独立分量, 
恰好可以组织为两组三维矢量。约定
\begin{subequations}\label{eq:EB-def}
\begin{align}
    E_i &\equiv f_{0i} \\
    B_k &\equiv -\tfrac{1}{2}\epsilon_{ijk}f_{ij}
\end{align}
\end{subequations}
即 $f_{ij} = -\epsilon_{ijk}B_k$, 
逆关系 $B_k = \frac{1}{2}\epsilon_{ijk}f_{ij}$ 
($B_1 = f_{32}$, $B_2 = f_{13}$, $B_3 = f_{21}$)。
将全部分量列出：
\begin{equation}
    f_{\mu\nu} = \begin{pmatrix}
    0 & E_1 & E_2 & E_3 \\
    -E_1 & 0 & -B_3 & B_2 \\
    -E_2 & B_3 & 0 & -B_1 \\
    -E_3 & -B_2 & B_1 & 0
    \end{pmatrix}
    \label{eq:f-matrix}
\end{equation}
定义 $f_{\mu\nu}$ 的 Hodge 对偶张量为
\begin{equation}
    (\star f)_{\mu\nu} \equiv \tfrac{1}{2}\,\epsilon_{\mu\nu\rho\sigma}\,f^{\rho\sigma}
    \label{eq:hodge-def}
\end{equation}
其中 $\epsilon_{0123}=-1$ ($\epsilon^{0123}=+1$)。
显式计算：
\begin{subequations}\label{eq:hodge-components}
\begin{align}
    (\star f)_{0i}
    &= \tfrac{1}{2}\,\epsilon_{0ijk}\,f^{jk} 
    = \tfrac{1}{2}\,\epsilon_{ijk}\,f^{jk} = B_i \\
    (\star f)_{ij}
    &= \tfrac{1}{2}\,\epsilon_{ij\rho\sigma}\,f^{\rho\sigma}
    = \epsilon_{ij0k}\,f^{0k}
    = -\epsilon_{ijk}\,f^{0k} = \epsilon_{ijk}\,E_k
\end{align}
\end{subequations}
对偶张量具有与 (\ref{eq:f-matrix}) 相同的结构, 
但 $\boldsymbol{E}$ 和 $\boldsymbol{B}$ 互换并带上符号：
\begin{equation}
    (\star f)_{\mu\nu}: \quad \boldsymbol{E} \to \boldsymbol{B}, \qquad 
    \boldsymbol{B} \to -\boldsymbol{E}
    \label{eq:hodge-EB}
\end{equation}
矩阵形式为
\begin{equation}
    (\star f)_{\mu\nu} = \begin{pmatrix}
    0 & B_1 & B_2 & B_3 \\
    -B_1 & 0 & E_3 & -E_2 \\
    -B_2 & -E_3 & 0 & E_1 \\
    -B_3 & E_2 & -E_1 & 0
    \end{pmatrix}
\end{equation}
Hodge 对偶的平方：再作用一次,
\begin{equation}
    \star^2 f_{\mu\nu} = \tfrac{1}{2}\epsilon_{\mu\nu\rho\sigma}
    \bigl(\tfrac{1}{2}\epsilon^{\rho\sigma\alpha\beta} f_{\alpha\beta}\bigr)
    = -f_{\mu\nu}
    \label{eq:hodge-squared}
\end{equation}
(闵氏时空中 Hodge 对偶平方为 $-1$)。
由于 $\star^2=-1$, 算符 $i\star$ 满足 $(i\star)^2=1$, 本征值为 $\pm 1$。
投影算符
\begin{equation}
    \mathbb{P}^\pm_{\mu\nu}{}^{\rho\sigma}
    = \tfrac{1}{2}\bigl(
    \delta^\rho_\mu\delta^\sigma_\nu - \delta^\rho_\nu\delta^\sigma_\mu
    \pm \tfrac{i}{2}\,\epsilon_{\mu\nu}{}^{\rho\sigma}\bigr)
\end{equation}
满足 $\mathbb{P}^++\mathbb{P}^-=\mathbb{I}$,
$\mathbb{P}^+\mathbb{P}^-=0$, $(\mathbb{P}^\pm)^2=\mathbb{P}^\pm$,
给出分解
\begin{subequations}\label{eq:f-pm-decomp}
\begin{align}
    f^\pm_{\mu\nu} &\equiv \mathbb{P}^\pm_{\mu\nu}{}^{\rho\sigma} f_{\rho\sigma}
    = \tfrac{1}{2}\bigl(f_{\mu\nu} \pm i(\star f)_{\mu\nu}\bigr) \\
    f_{\mu\nu} &= f^+_{\mu\nu} + f^-_{\mu\nu}, \qquad
    i\star f^\pm_{\mu\nu} = \pm f^\pm_{\mu\nu}
\end{align}
\end{subequations}
在三维矢量语言下, 该分解对应 Riemann--Silberstein 复矢量
\begin{equation}
    \boldsymbol{f}^{\,\pm} = \tfrac{1}{2}(\boldsymbol{E} \pm i\boldsymbol{B})
    \label{eq:RS-def}
\end{equation}
满足 $\boldsymbol{f}^{\,-}=(\boldsymbol{f}^{\,+})^*$。


\paragraph*{6 维表示空间的不可约分解。}
上面的分解 (\ref{eq:RS-def}) 是代数定义。
现在验证 $\boldsymbol{f}^{\,+}$ 确实按 $(0,1)$ 变换, 
$\boldsymbol{f}^{\,-}$ 按 $(1,0)$ 变换。

沿 $z$ 方向快度为 $\phi$ 的标准 boost 下,
$f'_{\mu\nu}=\Lambda^\rho_{\ \mu}\Lambda^\sigma_{\ \nu}f_{\rho\sigma}$ 给出
\begin{subequations}\label{eq:boost-EB}
\begin{align}
    E'_1 &= \cosh\phi\,E_1 - \sinh\phi\,B_2, &\quad 
    E'_2 &= \cosh\phi\,E_2 + \sinh\phi\,B_1, &\quad 
    E'_3 &= E_3 \label{eq:boost-E}\\
    B'_1 &= \sinh\phi\,E_2 + \cosh\phi\,B_1, &\quad  
    B'_2 &= -\sinh\phi\,E_1 + \cosh\phi\,B_2, &\quad  
    B'_3 &= B_3 \label{eq:boost-B}
\end{align}
\end{subequations}
boost 混合 $\boldsymbol{E}$ 和 $\boldsymbol{B}$, 
而旋转把 $\boldsymbol{E}$ 和 $\boldsymbol{B}$ 各自内部旋转——
6 维空间 $(\boldsymbol{E},\boldsymbol{B})$ 在 Lorentz 群下不是不可约的。
计算 $f'^+_1=\frac{1}{2}(E'_1+iB'_1)$：
\begin{align}
    f'^+_1 &= \tfrac{1}{2}\bigl[(\cosh\phi\,E_1-\sinh\phi\,B_2)
    +i(\sinh\phi\,E_2+\cosh\phi\,B_1)\bigr] \nonumber\\
    &= \tfrac{1}{2}\bigl[\cosh\phi\,(E_1+iB_1)+i\sinh\phi\,(E_2+iB_2)\bigr]
    = \cosh\phi\,f^+_1 + i\sinh\phi\,f^+_2
\end{align}
类似地
$f'^+_2 = i\sinh\phi\,f^+_1+\cosh\phi\,f^+_2$, $f'^+_3=f^+_3$。
$\boldsymbol{f}^{\,+}$ 的三个分量只和自己混合, 完全不涉及 $\boldsymbol{f}^{\,-}$;
$\boldsymbol{f}^{\,-}$ 同理 ($i\to-i$)。
6 维空间在 boost 下分裂为两个 3 维不变子空间。

旋转不区分 $\boldsymbol{E}$ 和 $\boldsymbol{B}$,
因此 $\boldsymbol{f}^{\,\pm}$ 的旋转生成元都是 $j=1$ 角动量矩阵 $J_i^{(1)}$,
在 Cartesian 基下 $(J^{(1)}_k)_{ij}=-i\epsilon_{kij}$。
但 Cartesian 基不是 $J_3$ 的本征态——
$J^{(1)}_3$ 的显式形式为
\begin{equation}
    J^{(1)}_3 = \begin{pmatrix} 0&-i&0\\i&0&0\\0&0&0 \end{pmatrix}
    \qquad\Longrightarrow\qquad
    J^{(1)}_3 \begin{pmatrix} f_1\\f_2\\f_3 \end{pmatrix}
    = \begin{pmatrix} -if_2\\if_1\\0 \end{pmatrix}
\end{equation}
为对角化 $J_3$, 引入球面基
\begin{subequations}\label{eq:spherical-basis}
\begin{align}
    f_+ &\equiv -\tfrac{1}{\sqrt{2}}(f_1+if_2), &\quad
    f_0 &\equiv f_3, &\quad
    f_- &\equiv \tfrac{1}{\sqrt{2}}(f_1-if_2)
    \label{eq:spherical-basis-def}\\[4pt]
    f_1 &= -\tfrac{1}{\sqrt{2}}(f_+-f_-), &\quad
    f_2 &= -\tfrac{i}{\sqrt{2}}(f_++f_-), &\quad
    f_3 &= f_0
    \label{eq:spherical-basis-inv}
\end{align}
\end{subequations}
直接验证
\begin{subequations}
    \begin{align}
    &J_3 f_+ = -\frac{1}{\sqrt{2}}(J_3 f_1+iJ_3 f_2)= -\frac{1}{\sqrt{2}}(-if_2-f_1) = +1\cdot f_+\\
    &J_3 f_0=0\\
    &J_3 f_-=-1\cdot f_-
\end{align}
\end{subequations}
也就是
\begin{equation}
    J_3\,f_m = m\,f_m, \qquad m=+1,\,0,\,-1
    \label{eq:J3-diag}
\end{equation}
将 boost 变换转到球面基, 利用 (\ref{eq:spherical-basis-def})：
\begin{subequations}\label{eq:boost-spherical}
\begin{align}
    f'^+_m &= e^{+m\phi}\,f^+_m \label{eq:fplus-boost}\\
    f'^-_m &= e^{-m\phi}\,f^-_m \label{eq:fminus-boost}
\end{align}
\end{subequations}
有限 boost $e^{-i\phi K_3}$ 在 $f^+_m$ 上给出 $e^{+m\phi}$,
即 $-iK_3^{(+)}=+J_3^{(1)}$\,\footnote{%
$J_3^{(1)}$ 在权基上 $J_3 f_m=mf_m$, 所以 $e^{+\phi J_3}f_m=e^{+m\phi}f_m$。};
推广到三个方向\footnote{%
对 $x,y$ 方向 boost 同理, 或由 $[J_i,K_j]=i\epsilon_{ijk}K_k$ 的代数一致性得到。}：
\begin{subequations}\label{eq:K-on-fpm}
\begin{align}
    \vec{K}^{(\boldsymbol{f}^+)} &= -i\boldsymbol{J}^{(1)} \\
    \vec{K}^{(\boldsymbol{f}^-)} &= +i\boldsymbol{J}^{(1)}
\end{align}
\end{subequations}
代入 $\boldsymbol{\mathcal{A}}=\frac{1}{2}(\boldsymbol{J}-i\boldsymbol{K})$,
$\boldsymbol{\mathcal{B}}=\frac{1}{2}(\boldsymbol{J}+i\boldsymbol{K})$：
\begin{subequations}\label{eq:AB-identification}
\begin{align}
    \boldsymbol{f}^+: &\quad 
    \mathcal{A}_i &&= \tfrac{1}{2}(J_i-i(-iJ_i)) &&= 0, &\quad
    \mathcal{B}_i &&= \tfrac{1}{2}(J_i+i(-iJ_i)) &&= J_i 
    &\quad\Longrightarrow\quad (0,1) \\
    \boldsymbol{f}^-: &\quad
    \mathcal{A}_i &&= \tfrac{1}{2}(J_i-i(+iJ_i)) &&= J_i, &\quad
    \mathcal{B}_i &&= \tfrac{1}{2}(J_i+i(+iJ_i)) &&= 0
    &\quad\Longrightarrow\quad (1,0)
\end{align}
\end{subequations}
与 $(A,B)$ 表示的 boost 矩阵 (3.64) 
$\mathscr{D}(L(p))=\exp(-\phi\cdot\boldsymbol{\mathcal{A}})\otimes\exp(+\phi\cdot\boldsymbol{\mathcal{B}})$ 
一致性检验：
\begin{subequations}\label{eq:boost-check}
\begin{align}
    (0,1)&: \; e^{+\phi B_3} \text{ 的本征值 } e^{+b\phi}\;(b=+1,0,-1)
    \;\longleftrightarrow\; f'^+_m=e^{+m\phi}f^+_m \;\checkmark \\
    (1,0)&: \; e^{-\phi A_3} \text{ 的本征值 } e^{-a\phi}\;(a=+1,0,-1)
    \;\longleftrightarrow\; f'^-_m=e^{-m\phi}f^-_m \;\checkmark
\end{align}
\end{subequations}
由此确认
\begin{equation}
    \boxed{
    \underbrace{f^+_{\mu\nu}}_{(0,1),\;\lambda=+1} \oplus 
    \underbrace{f^-_{\mu\nu}}_{(1,0),\;\lambda=-1} = f_{\mu\nu}
    }
    \label{eq:dictionary}
\end{equation}


\paragraph*{从球面基到极化矢量。}(\ref{eq:dictionary}) 建立后, 
ISO(2) 约束 (3.91) 要求物理态锁定为：
\begin{subequations}\label{eq:ISO2-select}
\begin{align}
    \lambda=+1,\;(0,1): &\quad a=0,\;b=+1 
    \quad\Longrightarrow\quad \text{仅 }f^+_+\text{ 存活} \\
    \lambda=-1,\;(1,0): &\quad a=-1,\;b=0 
    \quad\Longrightarrow\quad \text{仅 }f^-_-\text{ 存活}
\end{align}
\end{subequations}
在标准动量处, 全部场强信息被编码在唯一一个球面基分量中。
下面将球面基还原为 Cartesian 分量, 分别处理两种螺旋度。

$\lambda=+1$: 仅 $(0,1)$ 扇区的 $b=+1$ 非零, $(1,0)$ 扇区不贡献, 
即 $\boldsymbol{f}^-=0$。
由 (\ref{eq:spherical-basis-def}), 
$|b=+1\rangle = -\frac{1}{\sqrt{2}}(\hat{e}_1+i\hat{e}_2)$, 
这给出了 $\boldsymbol{f}^+$ 的 Cartesian 方向。
由逆关系 (\ref{eq:RS-def}):
\begin{equation}
    \boldsymbol{E} = \boldsymbol{f}^+ + \underbrace{\boldsymbol{f}^-}_{=0} 
    = \boldsymbol{f}^+, \qquad
    \boldsymbol{B} = -i(\boldsymbol{f}^+ - \underbrace{\boldsymbol{f}^-}_{=0}) 
    = -i\boldsymbol{f}^+ = -i\boldsymbol{E}
\end{equation}
因此 $\boldsymbol{E}$ 的方向就是 $\boldsymbol{f}^+$ 的方向：
\begin{equation}
    \boldsymbol{E} = \boldsymbol{f}^+ 
    \propto -\tfrac{1}{\sqrt{2}}(\hat{e}_1+i\hat{e}_2) 
    = \tfrac{1}{\sqrt{2}}(-1,-i,0)
    \label{eq:E-plus}
\end{equation}

$\lambda=-1$: 仅 $(1,0)$ 扇区的 $a=-1$ 非零, $(0,1)$ 扇区不贡献, 
即 $\boldsymbol{f}^+=0$。
由 (\ref{eq:spherical-basis-def}), 
$|a=-1\rangle = \frac{1}{\sqrt{2}}(\hat{e}_1-i\hat{e}_2)$, 
这给出了 $\boldsymbol{f}^-$ 的 Cartesian 方向。
由逆关系 (\ref{eq:RS-def}):
\begin{equation}
    \boldsymbol{E} = \underbrace{\boldsymbol{f}^+}_{=0} + \boldsymbol{f}^- 
    = \boldsymbol{f}^-, \qquad
    \boldsymbol{B} = -i(\underbrace{\boldsymbol{f}^+}_{=0} - \boldsymbol{f}^-) 
    = +i\boldsymbol{f}^- = +i\boldsymbol{E}
\end{equation}
因此
\begin{equation}
    \boldsymbol{E} = \boldsymbol{f}^- 
    \propto \tfrac{1}{\sqrt{2}}(\hat{e}_1-i\hat{e}_2)
    = \tfrac{1}{\sqrt{2}}(+1,-i,0)
    \label{eq:E-minus}
\end{equation}

两种螺旋度的结果可以对比写为
\begin{subequations}\label{eq:E-from-spherical}
\begin{align}
    \lambda=+1: &\quad \boldsymbol{E} \propto \tfrac{1}{\sqrt{2}}(-1,-i,0), 
    \qquad \boldsymbol{B} = -i\boldsymbol{E} \\
    \lambda=-1: &\quad \boldsymbol{E} \propto \tfrac{1}{\sqrt{2}}(+1,-i,0),
    \qquad \boldsymbol{B} = +i\boldsymbol{E}
\end{align}
\end{subequations}
$\lambda=+1$ 对应 $\boldsymbol{E}$ 和 $\boldsymbol{B}$ 右旋旋转 (球面基 $m=+1$), 
$\lambda=-1$ 对应左旋 ($m=-1$)——
这正是圆极化的定义。

\paragraph*{从三维极化方向到四维极化矢量。}
在标准动量 $k^\mu=(\kappa,0,0,\kappa)$ 处, 光子沿 $z$ 传播,
电场在横向平面 $xy$ 内振荡。
(\ref{eq:E-from-spherical}) 给出空间分量, 
纵向分量 $\epsilon^3=0$ (电场垂直于传播方向),
时间分量由横向性 $k_\mu\epsilon^\mu=0$ 确定：
\begin{equation}
    k_\mu\epsilon^\mu = \kappa\,\epsilon^0 - \kappa\,\underbrace{\epsilon^3}_{=0} = 0
    \quad\Longrightarrow\quad \epsilon^0=0
\end{equation}
四维极化矢量为
\begin{subequations}\label{eq:circ-pol}
\begin{align}
    \epsilon^\mu(+1) &= \tfrac{1}{\sqrt{2}}(0,-1,-i,0) \\
    \epsilon^\mu(-1) &= \tfrac{1}{\sqrt{2}}(0,+1,-i,0)
\end{align}
\end{subequations}
满足 $k\cdot\epsilon=0$, $\epsilon\cdot\epsilon^*=-1$,
$J_3\epsilon(\lambda)=\lambda\epsilon(\lambda)$。
圆极化矢量不是凭空约定的, 而是球面基 + ISO(2) 约束嵌入四维的唯一产物：
\begin{equation*}
    \underbrace{J_3\text{ 对角化}}_{\text{球面基}}
    \xrightarrow{\;\text{ISO(2)}\;}
    \underbrace{|m=\pm 1\rangle\text{ 存活}}_{\text{圆极化方向}}
    \xrightarrow{\;\text{嵌入四维}\;}
    \underbrace{\epsilon^\mu(\lambda)}_{\text{极化矢量}}
\end{equation*}


\paragraph*{构造反对称极化张量。}
用 $k_\mu$ 和 $\epsilon_\mu$ 构造反对称张量, 
约束为
\begin{enumerate}
    \item $e_{\mu\nu}=-e_{\nu\mu}$;
    \item $k^\mu e_{\mu\nu}=0$;
    \item $\lambda=+1$ 自对偶, $\lambda=-1$ 反自对偶。
\end{enumerate} 
最一般的反对称组合为\footnote{%
$k_\mu k_\nu$ 和 $\epsilon_\mu\epsilon_\nu$ 反对称化后为零, 
故已是最一般形式。两项均自动满足横向性。}
\begin{equation}
    e_{\mu\nu} = \alpha\,\underbrace{(k_\mu\epsilon_\nu-k_\nu\epsilon_\mu)}_{T^{(1)}_{\mu\nu}}
    +\beta\,\underbrace{\epsilon_{\mu\nu\rho\sigma}k^\rho\epsilon^\sigma}_{T^{(2)}_{\mu\nu}}
    \label{eq:general-ansatz}
\end{equation}
圆极化矢量已编码对偶性——
在标准动量处直接计算\footnote{%
以 $\lambda=+1$ 为例:
$T^{(2)}_{01}=\epsilon_{0132}k^3\epsilon^2(+1)
=(+1)\kappa\frac{-i}{\sqrt{2}}=\frac{-i\kappa}{\sqrt{2}}$;
$T^{(1)}_{01}=k_0\epsilon_1(+1)=\kappa\frac{1}{\sqrt{2}}$;
比值 $T^{(2)}_{01}/T^{(1)}_{01}=-i$。其余分量同。}验证：
\begin{equation}
    \epsilon_{\mu\nu\rho\sigma}k^\rho\epsilon^\sigma(\lambda)
    = -i\lambda(k_\mu\epsilon_\nu(\lambda)-k_\nu\epsilon_\mu(\lambda))
    \label{eq:duality-auto}
\end{equation}
代入 (\ref{eq:general-ansatz}):
$e_{\mu\nu}=(\alpha-i\lambda\beta)T^{(1)}$, 
两项合并为同一结构, 只剩一个自由常数。
选取 $\alpha=i$, $\beta=0$\,\footnote{%
$\beta\neq 0$ 由 (\ref{eq:duality-auto}) 不改变 $e_{\mu\nu}$ 的值。
$\alpha=i$ 使 $e_{0i}$ 直接给出电场极化 (差 $k^0$ 因子)。}:
\begin{equation}
    \boxed{e_{\mu\nu}(\mathbf{k},\lambda)
    = i\bigl(k_\mu\epsilon_\nu(\mathbf{k},\lambda)
    -k_\nu\epsilon_\mu(\mathbf{k},\lambda)\bigr)}
    \label{eq:pol-tensor}
\end{equation}
形式上与 $\partial_\mu A_\nu-\partial_\nu A_\mu$ 在动量空间相同,
但 $A_\mu$ 尚未引入——(\ref{eq:pol-tensor}) 完全来自
$(1,0)\oplus(0,1)$ 的表示论。

\paragraph*{验证自对偶性。}
以 $\lambda=+1$ 为例, 取 $k^\mu=(\kappa,0,0,\kappa)$,
$\epsilon_\mu=\eta_{\mu\nu}\epsilon^\nu=\frac{1}{\sqrt{2}}(0,1,i,0)$。
由 (\ref{eq:pol-tensor}) 计算\footnote{%
$k_3=\eta_{33}k^3=-\kappa$, 所以 $e_{3i}=-e_{0i}$
(即 $k^2=0$ 的推论)。}：
\begin{subequations}\label{eq:e-components}
\begin{align}
    e_{01} &= ik_0\epsilon_1 = \tfrac{i\kappa}{\sqrt{2}}, &\quad
    e_{02} &= ik_0\epsilon_2 = -\tfrac{\kappa}{\sqrt{2}} \\
    e_{31} &= -ik_3\epsilon_1 = -\tfrac{i\kappa}{\sqrt{2}}, &\quad
    e_{32} &= -ik_3\epsilon_2 = \tfrac{\kappa}{\sqrt{2}}
\end{align}
\end{subequations}
读出\footnote{%
$E_i=e_{0i}$; $B_1=e_{32}$, $B_2=e_{13}=-e_{31}$, $B_3=e_{21}=0$。}：
\begin{subequations}\label{eq:EB-from-pol}
\begin{align}
    \boldsymbol{E} &= \tfrac{\kappa}{\sqrt{2}}(i,-1,0) \\
    \boldsymbol{B} &= \tfrac{\kappa}{\sqrt{2}}(1,i,0)
\end{align}
\end{subequations}
验证：
\begin{subequations}\label{eq:self-dual-check}
\begin{align}
    \boldsymbol{E}+i\boldsymbol{B} &= \sqrt{2}\kappa(i,-1,0) \neq 0 \\
    \boldsymbol{E}-i\boldsymbol{B} &= \boldsymbol{0} \quad\checkmark
\end{align}
\end{subequations}
纯自对偶 $(0,1)$。$\lambda=-1$ 类似, 纯反自对偶 $(1,0)$。也就是说$\lambda=+1$ 的光子, 
$\boldsymbol{E}$ 和 $\boldsymbol{B}$ 在横向平面内右旋旋转,
$\boldsymbol{B}=-i\boldsymbol{E}$; 
$\lambda=-1$ 左旋, $\boldsymbol{B}=+i\boldsymbol{E}$。


\topictitle{场展开}
光子是自身的反粒子 ($a^{c\dagger}=a^\dagger$),
$(-1)^{2B}=1$ 对两个扇区均成立\footnote{%
$(1,0)$: $B=0$, $(-1)^{2B}=1$; $(0,1)$: $B=1$, $(-1)^{2B}=1$。};
由 (3.100), $u_{\mu\nu}=v_{\mu\nu}^*=e^*_{\mu\nu}$。场展开为
\begin{equation}
    \boxed{
    f_{\mu\nu}(x) = \frac{1}{(2\pi)^{3/2}} \sum_{\lambda=\pm 1} 
    \int \frac{d^3k}{\sqrt{2k^0}}
    \left[e_{\mu\nu}(\mathbf{k},\lambda)\,a(\mathbf{k},\lambda)\,e^{ik\cdot x}
    + e^*_{\mu\nu}(\mathbf{k},\lambda)\,a^\dagger(\mathbf{k},\lambda)\,e^{-ik\cdot x}\right]
    }
    \label{eq:f-expansion}
\end{equation}
其中 $k^0=|\mathbf{k}|$, 对易关系为
\begin{subequations}\label{eq:photon-ccr}
\begin{align}
    [a(\mathbf{k},\lambda),\,a^\dagger(\mathbf{k}',\lambda')] 
    &= \delta^3(\mathbf{k}-\mathbf{k}')\,\delta_{\lambda\lambda'} \\
    [a(\mathbf{k},\lambda),\,a(\mathbf{k}',\lambda')] 
    &= [a^\dagger(\mathbf{k},\lambda),\,a^\dagger(\mathbf{k}',\lambda')] = 0
\end{align}
\end{subequations}
逐项解读：
\begin{itemize}
    \item $a(\mathbf{k},\lambda)\,e^{ik\cdot x}$ (正频): 湮灭动量 $\mathbf{k}$、螺旋度 $\lambda$ 的光子;
    $a^\dagger(\mathbf{k},\lambda)\,e^{-ik\cdot x}$ (负频): 产生同一光子。
    \item $e_{\mu\nu}$, $e^*_{\mu\nu}$: 极化张量, 编码该模式的 $\boldsymbol{E}$-$\boldsymbol{B}$ 构型。
    \item $1/\sqrt{2k^0}$: Lorentz 不变归一化, 使 $d^3k/(2k^0)$ 构成不变测度。
\end{itemize}
正频与负频的同时出现是微观因果性的要求 (3.1 节 (3.22))。
由于光子是自身的反粒子, 不需引入额外的反粒子算符。

\textbf{厄米性。}对 (\ref{eq:f-expansion}) 取厄米共轭,
$a\leftrightarrow a^\dagger$, $e_{\mu\nu}\leftrightarrow e^*_{\mu\nu}$,
$e^{ikx}\leftrightarrow e^{-ikx}$, 恰好交换两项, 因此
$f_{\mu\nu}^\dagger(x) = f_{\mu\nu}(x)$。
经典极限中, 这对应 $\boldsymbol{E}$, $\boldsymbol{B}$ 是实矢量场。

\textbf{经典极限。}对相干态 $|\alpha\rangle$
($a(\mathbf{k},\lambda)|\alpha\rangle = \alpha(\mathbf{k},\lambda)|\alpha\rangle$) 取期望值,
得到经典场强张量的 Fourier 展开。对于单色平面波, 期望值化为
$\langle f_{\mu\nu}\rangle \propto e_{\mu\nu}(\mathbf{k}_0,\lambda)\,e^{ik_0\cdot x} + \text{c.c.}$
——传播方向 $\hat{k}_0$、圆极化 $\lambda$ 的经典平面电磁波, 极化张量的命名由此而来。

\textbf{极化求和。}利用 $e_{\mu\nu}=i(k_\mu\epsilon_\nu - k_\nu\epsilon_\mu)$ 展开：
\begin{equation}
    \sum_{\lambda=\pm 1} e_{\mu\nu}(\mathbf{k},\lambda)\,e^*_{\rho\sigma}(\mathbf{k},\lambda)
    = k_\mu k_\rho\,P_{\nu\sigma} - k_\mu k_\sigma\,P_{\nu\rho}
    - k_\nu k_\rho\,P_{\mu\sigma} + k_\nu k_\sigma\,P_{\mu\rho}
\end{equation}
其中 $P_{\mu\nu}=\sum_\lambda \epsilon_\mu\epsilon^*_\nu$。
对于耦合到守恒流 $j^\mu$ ($k_\mu j^\mu=0$) 的过程,
Ward 恒等式允许 $P_{\mu\nu}\to -\eta_{\mu\nu}$。

\textbf{推广到非阿贝尔规范场。}
上述构造对\textbf{任何}零质量自旋-1 粒子都成立, 区别仅在于内部对称性。
对于规范群 $G$ 的规范玻色子, 场强张量携带伴随表示指标 $a=1,\ldots,\dim G$,
极化张量 $e_{\mu\nu}$ 与色指标无关——Lorentz 结构和内部对称性完全解耦。
对易关系推广为
$[a^a(\mathbf{k},\lambda),\,a^{b\dagger}(\mathbf{k}',\lambda')]
= \delta^{ab}\,\delta^3(\mathbf{k}-\mathbf{k}')\,\delta_{\lambda\lambda'}$。
在完整的非阿贝尔理论中, 场强张量包含非线性项
$F^a_{\mu\nu} = \partial_\mu A^a_\nu - \partial_\nu A^a_\mu + gf^{abc}A^b_\mu A^c_\nu$,
自由场展开仅对应其线性化部分。

这个场展开\textbf{不包含任何规范冗余}——
它直接描述了光子 (或胶子) 的两个物理自由度 $\lambda=\pm 1$。

\topictitle{矢量势与规范对称性}
描述光子的另一种方式是使用四矢量场 $A^\mu$, 按 $(\frac{1}{2},\frac{1}{2})$ 表示变换。
角动量加法 $\frac{1}{2}\otimes\frac{1}{2}=0\oplus 1$ 包含 $j=0$ 和 $j=1$ 两个扇区,
共 4 个分量, 但光子只有 $\lambda=\pm 1$ 两个物理自由度,
多出的分量正是规范自由度的来源。
理解规范冗余的关键在于 Wigner 小群 $ISO(2)$。
``平移"生成元 $\Pi_1,\Pi_2$ 在物理态上本征值为零 (排除连续自旋),
但在四矢量表示中, 小群变换作用在 $\epsilon^\mu$ 上时产生额外项：
\begin{equation}
    W^\mu_{\ \nu}\,\epsilon^\nu(\mathbf{k},\lambda) =
    e^{i\lambda\theta}\,\epsilon^\mu(\mathbf{k},\lambda) + \Omega(\theta,\alpha,\beta)\,k^\mu
\end{equation}
第二项 $\propto k^\mu$ 在坐标空间中对应 $A_\mu(x)\to A_\mu(x)+\partial_\mu\Omega(x)$
——这正是规范变换。

\begin{center}
\fbox{\parbox{0.88\textwidth}{
规范对称性不是额外施加的, 而是 Poincar\'{e} 群在零质量矢量场上表示论的自然推论。
$ISO(2)$ 小群的``平移"生成元在物理态上本征值为零,
但在 $A^\mu$ 上产生了 $\delta A_\mu = \partial_\mu\Omega$ 的规范变换。
}}
\end{center}

从表示论看, 取反对称导数实现了映射
\begin{equation}
    \underbrace{A^\mu}_{(\frac{1}{2},\frac{1}{2})}
    \xrightarrow{\;\partial_\mu A_\nu - \partial_\nu A_\mu\;}
    \underbrace{f_{\mu\nu}}_{(1,0)\oplus(0,1)}
\end{equation}
``投射掉"了 $j=0$ 分量, 只保留 $j=1$ 的两个横向自由度。

矢量势的场展开 (辐射规范 $A^0=0$, $\nabla\cdot\boldsymbol{A}=0$) 为
\begin{equation}
    A^\mu(x) = \frac{1}{(2\pi)^{3/2}} \sum_{\lambda=\pm 1}
    \int \frac{d^3k}{\sqrt{2k^0}}
    \left[\epsilon^\mu(\mathbf{k},\lambda)\,a(\mathbf{k},\lambda)\,e^{ik\cdot x}
    + \epsilon^{*\mu}(\mathbf{k},\lambda)\,a^\dagger(\mathbf{k},\lambda)\,e^{-ik\cdot x}\right]
\end{equation}
其中产生湮灭算符与 (\ref{eq:f-expansion}) 中的完全相同。
取反对称导数, 利用 $\partial_\mu e^{\pm ik\cdot x}=\pm ik_\mu e^{\pm ik\cdot x}$,
精确回到 (\ref{eq:f-expansion}): $\partial_\mu A_\nu - \partial_\nu A_\mu = f_{\mu\nu}$

\topictitle{为什么需要矢量势}

$f_{\mu\nu}$ 完整描述自由光子, 引入 $A^\mu$ 的理由在于相互作用：

\begin{enumerate}
    \item 局域相互作用:
    最小耦合 $\mathscr{H}_{\text{int}}=j^\mu A_\mu$ (维度 4, 可重整);
    $f_{\mu\nu}$ 耦合为 $j^{\mu\nu}f_{\mu\nu}$ (维度 5, 不可重整)。
    
    \item 规范不变性作为代价:
    $A^\mu$ 带来规范冗余, 需规范固定与 Faddeev-Popov 程序,
    换来可重整性与局域性。
    
    \item Aharonov-Bohm 效应:
    $\exp(ie\oint A_\mu dx^\mu)$ 在 $f_{\mu\nu}=0$ 区域仍可观测,
    $A_\mu$ 在量子理论中有超越 $f_{\mu\nu}$ 的物理内容。
\end{enumerate}

$(1,0)\oplus(0,1)$ 提供自由光子的完整描述;
$(\frac{1}{2},\frac{1}{2})$ 是引入相互作用的必要工具。


